\chapter{Công nợ --- Đổi trả · Bù hàng}

Chương này gồm \textbf{8 đường dẫn} gom thành \textbf{7 màn} cộng một đường dẫn cũ, thuộc hai
mô-đun: \rt{wujia_portal_debt} (3) và \rt{wujia_portal_return} (5).

\textbf{Bốn điểm cần biết trước khi đọc:}

\begin{enumerate}[leftmargin=*]
  \item \textbf{Hai phân hệ trong cùng một chương hiểu ``cửa hàng'' ngược nhau.} Công nợ chỉ
        làm việc với \textbf{đúng một cửa hàng đang chọn}; chưa chọn thì màn rỗng. Đổi trả thì
        \textbf{gộp mọi cửa hàng} tài khoản truy cập được --- một tài khoản ba cửa hàng mở màn
        Đổi trả sẽ thấy lẫn phiếu của cả ba. Cùng với ba kiểu đã nêu ở chương trước, portal
        hiện có \textbf{bốn cách} xác định ``cửa hàng đang xem''.
  \item \textbf{Công nợ là phân hệ chặn theo vai trò chặt nhất.} Nhân viên \textbf{không vào
        được cả ba màn} công nợ --- không phải ẩn số liệu, mà là một trang ``không có quyền''
        riêng. Và vai trò được xét tại \textbf{đúng cửa hàng đang chọn}, khác với Báo cáo
        (chương 5) xét vai trò cao nhất trên mọi cửa hàng.
  \item \textbf{Trên UAT chỉ 1 trong 4 cửa hàng có chứng từ kế toán.} Toàn bộ \uat{6} chứng từ
        đã ghi sổ đều thuộc \emph{TP HCM Quận 1}. Ba cửa hàng còn lại mở màn Công nợ lúc nào
        cũng thấy trạng thái rỗng --- đó là do dữ liệu, không phải do luật lọc sai.
  \item \textbf{Phiếu đổi trả lưu nháp là ngõ cụt.} Màn tạo yêu cầu trên máy tính có nút
        \emph{Lưu nháp}, nhưng portal \textbf{không có đường nào} để gửi một phiếu nháp đã lưu:
        màn chi tiết chỉ đọc, không có nút gửi. UAT đang có \uat{3} phiếu nháp nằm vĩnh viễn ở
        đó. Xem mục 7 của màn Tạo yêu cầu.
\end{enumerate}

Quy ước chung toàn portal (phân trang, định dạng ngày giờ và tiền, huy hiệu trạng thái, vai
trò, đính kèm) đã viết một lần ở mục~\ref{sec:quy-uoc-chung}; chương này chỉ nói phần khác
biệt.

% =============================================================================
\section{Màn Công nợ theo tuần}
\rt{/portal/debt}
\medskip

\mvao

Mục \emph{Công nợ} trên thanh điều hướng bên trái (máy tính) hoặc dòng \emph{Công nợ \& thanh
toán} trong bảng \emph{Thêm} (điện thoại). Màn nhận các tham số trên thanh địa chỉ:
\rt{week} (tuần đang xem), \rt{all} (điện thoại: bung hết hoá đơn), \rt{page} ·
\rt{page_size} (phân trang bảng máy tính), \rt{notice} (băng thông báo khi vừa bị đẩy về từ
màn chuyển khoản).

\mai

\begin{itemize}[leftmargin=*]
  \item Phải đăng nhập.
  \item \textbf{Chỉ Chủ tiệm và Quản lý}. Nhân viên mở màn này sẽ nhận một \textbf{trang thông
        báo không có quyền} --- không phải lỗi hệ thống, và không lộ bất kỳ con số nào.
  \item Vai trò được xét tại \textbf{đúng cửa hàng đang chọn}. Một người vừa là Quản lý cửa
        hàng A vừa là Nhân viên cửa hàng B: đang chọn A thì vào được, bấm đổi sang B thì chính
        tài khoản đó bị chặn. Trên UAT có đúng \uat{1} tài khoản như vậy (Phạm Thị Dung ---
        Quản lý \emph{Đống Đa}, Nhân viên \emph{TP HCM Quận 1} và \emph{Cầu Giấy}).
  \item \textbf{Chưa chọn cửa hàng thì không bị chặn.} Không có cửa hàng thì không có vai trò
        để xét, nên hệ thống cho vào và hiện trạng thái rỗng ``chưa chọn cửa hàng''. Nghĩa là
        một Nhân viên nhiều cửa hàng chưa bấm chọn vẫn mở được màn --- chỉ là không thấy số.
  \item Số liệu đọc bằng quyền hệ thống; giới hạn duy nhất là \rt{franchise_id} bằng đúng cửa
        hàng đang chọn.
  \item Thực tế UAT: \uat{4} cửa hàng, \uat{10} lượt tham gia, trong đó \uat{6} lượt đạt Chủ
        tiệm/Quản lý. Nhưng chỉ \emph{TP HCM Quận 1} có chứng từ, nên rốt cuộc chỉ \uat{2} tài
        khoản thực sự nhìn thấy số: chủ cửa hàng đó và Administrator (Quản lý cùng cửa hàng).
\end{itemize}

\mhien

\begin{hienbang}
Ô chọn tuần &
  --- (máy tự sinh) &
  \textbf{6 tuần}: tuần hiện tại và 5 tuần trước, tính từ thứ Hai. Nhãn dạng
  \emph{``10/08 -- 16/08/2026''}. Tuần cũ hơn \textbf{không chọn tới được}. &
  Mới nhất trước &
  Tại ngày chụp: từ 14/09 lùi về 10/08/2026 \\ \addlinespace
Tuần mở sẵn &
  \rt{account.move} &
  Thứ tự ưu tiên: (1) tuần \textbf{quá hạn cũ nhất còn dư nợ}; (2) nếu không có, tuần liền
  trước nếu tuần đó có chứng từ; (3) nếu vẫn không, tuần hiện tại. &
  --- &
  TP HCM Q1 mở sẵn tuần \uat{10/08 -- 16/08} \\ \addlinespace
Thẻ tổng của tuần &
  \rt{account.move} &
  Hoá đơn bán và giấy báo có \textbf{đã ghi sổ} (\rt{state = posted},
  \rt{move_type in (out_invoice, out_refund)}) của cửa hàng, có \textbf{ngày hoá đơn} rơi vào
  thứ Hai--Chủ nhật của tuần. Giấy báo có \textbf{trừ} vào tổng. Chứng từ nháp không tính. &
  Ngày hoá đơn, rồi số thứ tự &
  \uat{6} chứng từ ghi sổ / \uat{6} chứng từ còn nháp \\ \addlinespace
Số \emph{Còn phải trả} &
  \rt{account.move.amount_residual} &
  Tổng số dư ròng, \textbf{cắt sàn ở 0}: âm thì đưa sang ô \emph{Dư có}. Không bao giờ in số
  nợ âm. &
  --- &
  Tuần 10/08: còn phải trả \uat{0} · dư có \uat{\$72.448,85} \\ \addlinespace
Trạng thái tuần &
  suy từ số dư &
  Năm nhãn, xét theo thứ tự: hết dư $\rightarrow$ \emph{Đã thanh toán}; dư âm $\rightarrow$
  \emph{Dư có}; có hoá đơn quá hạn $\rightarrow$ \emph{Có quá hạn}; trả được một phần
  $\rightarrow$ \emph{Thanh toán một phần}; còn lại \emph{Chưa thanh toán}. &
  --- &
  Tuần 10/08 hiện \emph{Dư có} \\ \addlinespace
Hạn thanh toán &
  máy tự tính &
  Thứ Hai của tuần \textbf{cộng 10 ngày} --- tức thứ Năm tuần kế. Không đọc hạn ghi trên hoá
  đơn. &
  --- &
  Tuần 10/08 $\Rightarrow$ hạn \uat{20/08/2026} \\ \addlinespace
Bảng hoá đơn (máy tính) &
  \rt{account.move} &
  Mỗi dòng: số chứng từ, ngày, hạn, trạng thái, \textbf{tổng}, \textbf{đã trả}, \textbf{còn
  phải trả}. Giấy báo có in số âm. &
  \textbf{10} dòng mỗi trang &
  Tuần 10/08 có \uat{5} chứng từ \\ \addlinespace
Khối hoá đơn (điện thoại) &
  \rt{account.move} &
  Chỉ \textbf{2 hoá đơn} đầu kèm dòng ``còn \emph{n} hoá đơn''; bung hết bằng \rt{?all=1}. &
  --- & --- \\ \addlinespace
Trạng thái từng hoá đơn &
  suy từ số dư &
  \emph{Đã thanh toán} (hết dư, hoặc kế toán đã đánh dấu tất toán) · \emph{Giấy báo có} ·
  \emph{Quá hạn} · \emph{Một phần} (đã trả được ít nhất một đồng) · \emph{Chưa thanh toán}. &
  --- &
  Tuần 10/08: \uat{3} Đã thanh toán · \uat{1} Giấy báo có · \uat{1} Quá hạn \\ \addlinespace
Hộp chuyển khoản &
  \rt{res.partner.bank} &
  Chỉ dựng khi \textbf{còn phải trả lớn hơn 0}. Lấy tài khoản của công ty có bật cờ nhận tiền
  qua portal, số thứ tự nhỏ nhất. &
  --- &
  \uat{1}/1 tài khoản đã bật cờ \\
\end{hienbang}

\textbf{Về loại tiền:} ký hiệu tiền lấy theo \textbf{loại tiền của chính các chứng từ trong
tuần}; tuần trộn nhiều loại tiền thì lùi về tiền công ty. Trên UAT toàn bộ chứng từ ghi bằng
USD nên màn hiển thị \$ --- đúng, khác với màn Báo cáo (chương 5) đang in ký hiệu đồng Việt Nam cho số
USD.

\mloc

\begin{itemize}[leftmargin=*]
  \item \textbf{Tuần} --- chọn trong đúng 6 khoá tuần kể trên. Gõ khoá lạ trên thanh địa chỉ
        thì \textbf{không báo lỗi}, hệ thống lặng lẽ quay về tuần mặc định.
  \item \textbf{Bung hết hoá đơn} (điện thoại) --- \rt{?all=1}; giá trị khác coi như không bung.
  \item \textbf{Số dòng mỗi trang} (máy tính) --- chỉ nhận \textbf{10, 20, 50}.
  \item \textbf{Băng thông báo} --- chỉ chấp nhận đúng một giá trị \rt{no_due}; mọi chuỗi khác
        bị bỏ qua, nên không ai nhét được câu chữ tuỳ ý vào trang qua thanh địa chỉ.
\end{itemize}

\mkiem{Thanh toán}

\begin{kiembang}
1 & Đã chọn được cửa hàng & Màn rỗng ``chưa chọn cửa hàng'', nút không hiện \\
2 & Vai trò tại cửa hàng đó là Chủ tiệm hoặc Quản lý & Trang ``không có quyền'' \\
3 & Tuần đang xem \textbf{còn phải trả lớn hơn 0} &
    Bị đẩy về chính màn này kèm băng báo không còn khoản phải trả --- \textbf{không} in số âm,
    \textbf{không} dựng mã QR \\
4 & Công ty đã khai tài khoản nhận tiền bật cho portal &
    Màn chuyển khoản hiện thông báo chưa cấu hình, không hiện ô số tài khoản trống \\
\end{kiembang}

\mghi

\textbf{Không ghi gì.} Xem số tài khoản hay mở mã QR \textbf{không tạo phiếu thanh toán}, không
đánh dấu hoá đơn. Nội dung chuyển khoản do máy chủ tự sinh từ mã cửa hàng, số tuần và số tiền
--- \textbf{không tin} số tiền trình duyệt gửi lên.

\mhoi

\begin{ask}
\begin{enumerate}[leftmargin=*]
  \item \textbf{Số công nợ ở Trang chủ và số ở màn này là hai phép tính khác nhau.} Thẻ Trang
        chủ và huy hiệu đọc \textbf{ô tổng lưu sẵn} trên cửa hàng: tính bằng \textbf{tiền công
        ty}, \textbf{không giới hạn tuần}, cập nhật bằng tác vụ nền hằng ngày. Màn này tính
        trực tiếp bằng \textbf{tiền của chứng từ} và \textbf{chỉ trong một tuần}. UAT chưa khai
        tỷ giá nên hai số đang trùng nhau; khi khai tỷ giá thật, hai chỗ sẽ ra hai con số khác
        nhau cho cùng một cửa hàng.
  \item \textbf{Dropdown chỉ 6 tuần.} UAT có \uat{1} hoá đơn ngày 05/08/2026 nằm ngoài khung
        --- cửa hàng \textbf{không có cách nào} xem lại. Anh muốn giữ 6 tuần, hay thêm ô chọn
        khoảng ngày tự do?
  \item \textbf{``Dư có'' đè lên ``Có quá hạn''.} Tuần 10/08 của TP HCM Q1 có một hoá đơn quá
        hạn \$1,15 nhưng thẻ tuần hiện \emph{Dư có} vì giấy báo có lớn hơn, và nút Thanh toán
        biến mất. Cửa hàng nhìn vào tưởng không còn gì phải xử lý. Anh chốt giúp: quá hạn có
        được ưu tiên hiện không?
  \item \textbf{Trạng thái rỗng nói ``không còn công nợ''.} Câu này dựa trên ô tổng lưu sẵn có
        lớn hơn 0 hay không. TP HCM Q1 đang có ô tổng \textbf{âm}, nên mọi tuần không có chứng
        từ đều báo ``không còn công nợ'' --- trong khi vẫn còn một hoá đơn quá hạn.
  \item \textbf{Hạn thanh toán do portal tự tính} (thứ Hai + 10 ngày), \textbf{không} đọc hạn
        ghi trên hoá đơn. UAT có hoá đơn hạn 12/08 trong khi portal nói hạn tuần là 20/08.
        Lấy theo cái nào?
  \item \textbf{Tuần mở sẵn} = tuần quá hạn cũ nhất còn dư. Anh xác nhận cách mở này.
  \item \textbf{Ảnh mã QR vẫn đang treo} --- hiện chỉ có số tài khoản và nội dung chuyển khoản.
        Anh chốt giúp: mã QR tĩnh dán sẵn hay mã động sinh theo số tiền từng tuần?
  \item \textbf{Nhân viên bị chặn khỏi toàn bộ phân hệ Công nợ.} Xác nhận đúng ý anh.
\end{enumerate}
\end{ask}

% =============================================================================
\section{Màn Lịch sử thanh toán}
\rt{/portal/debt/payment-history}
\medskip

\mvao

Thẻ \emph{Lịch sử thanh toán} ngay trong phân hệ Công nợ. Tham số: \rt{month} (kỳ),
\rt{q} (tìm kiếm), \rt{page} · \rt{page_size}. Ngoài ra còn nhận \rt{date_from} ·
\rt{date_to} nhưng giao diện \textbf{chưa có ô nào} sinh ra hai tham số này.

\mai

Giống hệt màn Công nợ: chỉ Chủ tiệm / Quản lý của \textbf{đúng cửa hàng đang chọn}; Nhân viên
nhận trang không có quyền; chưa chọn cửa hàng thì vào được nhưng rỗng.

\mhien

\begin{hienbang}
Ô chọn kỳ &
  --- (máy tự sinh) &
  \textbf{6 kỳ}, mỗi kỳ là \textbf{trọn một tháng}: tháng hiện tại và 5 tháng trước. Nhãn
  \emph{``01/09 -- 30/09/2026''}. &
  Mới nhất trước &
  Tại ngày chụp: 09/2026 lùi về 04/2026 \\ \addlinespace
Bảng giao dịch &
  \rt{account.payment} &
  Phiếu thanh toán của cửa hàng \textbf{đã được xác nhận}
  (\rt{state in (in_process, paid)}), có \textbf{ngày thanh toán} nằm trong kỳ. Phiếu nháp,
  huỷ, từ chối \textbf{không hiện}. &
  Ngày mới nhất trước; \textbf{10} dòng mỗi trang &
  \uat{5}/6 phiếu đã xác nhận \\ \addlinespace
Cột \emph{Số tiền} &
  \rt{account.payment.amount} &
  Tiền cửa hàng trả vào hiện số dương; \textbf{tiền Ngô Gia chi ngược trả lại} hiện
  \textbf{số âm}, để không bị cộng nhầm vào phần đã trả. Ký hiệu tiền lấy theo
  \textbf{từng giao dịch}. &
  --- &
  UAT chưa có giao dịch chi ngược \\ \addlinespace
Cột \emph{Giờ} &
  \rt{create_date} &
  Ngày thanh toán trong kế toán \textbf{không có giờ}, nên giờ hiển thị lấy từ
  \textbf{lúc phiếu được tạo}, đổi theo múi giờ tài khoản. Hai mốc này có thể lệch nhau. &
  --- & --- \\ \addlinespace
Khối tổng kỳ &
  \rt{account.payment} &
  Tổng \textbf{toàn kỳ} (không phải tổng trang đang xem), \textbf{tách riêng từng loại tiền}
  --- \textbf{không quy đổi tỷ giá}, không cộng gộp số khác loại tiền. &
  --- &
  Kỳ 08/2026 ra \textbf{2 dòng}: \uat{4} giao dịch USD + \uat{1} giao dịch VND \\
\end{hienbang}

\mloc

\begin{itemize}[leftmargin=*]
  \item \textbf{Kỳ} --- chọn trong 6 kỳ; khoá lạ thì lặng lẽ về tháng hiện tại.
  \item \textbf{Tìm kiếm} --- khớp một phần theo \textbf{mã phiếu} hoặc \textbf{nội dung
        thanh toán}. Không tìm theo số tiền, không tìm theo số hoá đơn được gán.
  \item \textbf{Khoảng ngày} (gọi thẳng qua thanh địa chỉ) --- nếu từ ngày lớn hơn đến ngày thì
        hệ thống \textbf{tự hoán đổi hai đầu} rồi chạy tiếp, \textbf{không báo lỗi}. Khác hẳn
        năm màn còn lại của portal (Lịch sử, Giao hàng, Báo cáo, Đổi trả) --- những màn đó dừng
        lại và báo tại thanh lọc.
  \item \textbf{Số dòng mỗi trang} --- 10, 20, 50.
\end{itemize}

\mkiem{một dòng giao dịch}

Không có hành động nào --- bảng chỉ đọc, không mở màn chi tiết.

\mghi

\textbf{Không ghi gì.}

\mhoi

\begin{ask}
\begin{enumerate}[leftmargin=*]
  \item \textbf{Kỳ mở sẵn là tháng hiện tại, và trên UAT tháng đó rỗng.} Toàn bộ \uat{5} giao
        dịch nằm ở tháng 08/2026, còn tháng 09 có \uat{0} $\Rightarrow$ cửa hàng mở màn ra thấy
        trắng, phải tự nghĩ đến việc đổi kỳ. Anh muốn mặc định mở \textbf{kỳ gần nhất có giao
        dịch} không?
  \item \textbf{Ngày ngược thì tự hoán đổi} thay vì báo lỗi --- lệch với bốn màn lọc khác.
        Thống nhất theo cách nào?
  \item \textbf{Tổng tách theo loại tiền, không quy đổi.} Kỳ 08/2026 trên UAT hiện hai dòng
        tổng. Anh xác nhận đây là cách trình bày muốn giữ.
  \item \uat{1}/6 phiếu thanh toán trên UAT \textbf{chưa gắn cửa hàng} nên không cửa hàng nào
        thấy. Cần chặn ở khâu kế toán không?
\end{enumerate}
\end{ask}

% =============================================================================
\section{Màn Chuyển khoản}
\rt{/portal/debt/pay}
\medskip

\mvao

Nút \emph{Thanh toán} ở màn Công nợ. Nhận tham số \rt{week} để biết đang trả cho tuần nào.

\mai

Như hai màn trên: chỉ Chủ tiệm / Quản lý của cửa hàng đang chọn.

\mhien

\begin{hienbang}
Khối số tiền &
  \rt{account.move} &
  Số \emph{còn phải trả} của tuần đang chọn --- tính lại tại chỗ, không tin số trình duyệt
  gửi lên &
  --- &
  Tuần mặc định của TP HCM Q1 đang bằng \uat{0} \\ \addlinespace
Khối tài khoản &
  \rt{res.partner.bank} &
  Tài khoản của công ty có bật cờ nhận tiền qua portal, số thứ tự nhỏ nhất: tên ngân hàng,
  chủ tài khoản, số tài khoản &
  --- &
  \uat{1} tài khoản (Vietcombank -- NGO GIA UAT DEBT) \\ \addlinespace
Nội dung chuyển khoản &
  máy tự sinh &
  \textbf{Mã cửa hàng + K + số tuần + số tiền} (làm tròn xuống, không bao giờ âm). Ví dụ
  tuần 10/08 là tuần ISO 33 $\Rightarrow$ \emph{HCM-01 K33 ...} &
  --- & --- \\ \addlinespace
Ảnh mã QR &
  \textbf{chưa có} &
  Đang chờ BA chốt loại mã (tĩnh hay động). Màn hiện số tài khoản và nội dung, không có ảnh. &
  --- & --- \\
\end{hienbang}

\mloc

Không có ô nhập nào. Chỉ tham số \rt{week}; khoá lạ thì về tuần mặc định.

\mkiem{vào màn này}

\begin{kiembang}
1 & Đã chọn được cửa hàng & Trạng thái rỗng \\
2 & Vai trò Chủ tiệm / Quản lý tại cửa hàng đó & Trang ``không có quyền'' \\
3 & Tuần còn phải trả lớn hơn 0 &
    Quay về màn Công nợ đúng tuần đó kèm băng báo không còn khoản phải trả.
    \textbf{Trên UAT đây là nhánh luôn xảy ra}: tuần mặc định đang dư có \\
4 & Có tài khoản ngân hàng bật portal & Màn báo chưa cấu hình thay vì hiện ô trống \\
\end{kiembang}

\mghi

\textbf{Không ghi gì.} Đây là màn hướng dẫn chuyển khoản, không phải cổng thanh toán --- hệ
thống không tạo phiếu thu, không đánh dấu hoá đơn, không gọi sang ngân hàng.

\mhoi

\begin{ask}
\begin{enumerate}[leftmargin=*]
  \item \textbf{Ảnh mã QR} --- điểm treo lâu nhất của phân hệ. Mã tĩnh (dán sẵn một ảnh cho mọi
        cửa hàng) hay mã động (sinh theo số tiền và nội dung từng tuần)?
  \item \textbf{Nội dung chuyển khoản làm tròn xuống} về số nguyên. Với tiền USD có phần lẻ
        (UAT đang có hoá đơn \$1,15) thì nội dung sẽ ghi \emph{1} --- lệch số thật. Cần giữ
        phần lẻ không?
  \item \textbf{Sau khi cửa hàng chuyển khoản thì portal không biết gì} --- phải đợi kế toán
        Ngô Gia ghi nhận rồi tác vụ nền cập nhật lại. Cửa hàng có cần chỗ báo ``tôi đã chuyển''
        kèm ảnh chụp màn hình không?
\end{enumerate}
\end{ask}

% =============================================================================
\section{Màn Danh sách yêu cầu đổi trả · bù hàng}
\rt{/portal/return}
\medskip

\mvao

Dòng \emph{Đổi trả / Bù hàng} trong bảng \emph{Thêm} (điện thoại); trên Trang chủ có một thẻ số
liệu và một khối danh sách cùng trỏ vào đây. Tiêu đề màn ghi là \emph{Bù hàng}.

\mai

\begin{itemize}[leftmargin=*]
  \item Phải đăng nhập. \textbf{Không chặn vai trò} --- Nhân viên xem và tạo yêu cầu bình
        thường, ngang Chủ tiệm.
  \item Phạm vi là \textbf{mọi cửa hàng tài khoản truy cập được}, \textbf{không} phải cửa hàng
        đang chọn. Người phụ trách ba cửa hàng thấy phiếu của cả ba trộn chung một danh sách,
        và không có bộ lọc theo cửa hàng.
  \item Không truy cập được cửa hàng nào thì màn hiện trạng thái rỗng.
  \item Dữ liệu đọc bằng quyền hệ thống. Hệ thống \textbf{có} khai một luật giới hạn cho người
        dùng portal (chỉ thấy phiếu của các cửa hàng mình tham gia) nhưng luật đó không được
        dùng tới ở màn này; điều kiện lọc trong bảng dưới cho ra đúng cùng kết quả.
\end{itemize}

\mhien

\begin{hienbang}
Danh sách phiếu &
  \rt{wujia.return.request} &
  Phiếu thuộc \textbf{mọi cửa hàng của tôi} (\rt{franchise_id in ...}). Không loại trạng
  thái nào --- kể cả nháp, từ chối, đã huỷ. &
  Ngày tạo mới nhất trước; \textbf{20} phiếu mỗi trang &
  \uat{14} phiếu: Cầu Giấy \uat{5} · Đống Đa \uat{2} · TP HCM Q1 \uat{7} \\ \addlinespace
Huy hiệu trạng thái &
  \rt{state} + tiến độ bù &
  Tám nhãn, trong đó hai nhãn là \textbf{ghép}: \emph{Đang xử lý} gộp hai trạng thái
  (\emph{đang xét} và \emph{đang thực hiện}); \emph{Đang bù một phần} \textbf{không phải một
  trạng thái} mà là ``đang thực hiện + đã bù được một phần''. &
  --- &
  \uat{3} Nháp · \uat{2} Đã gửi · \uat{5} Đã duyệt · \uat{2} Hoàn tất · \uat{2} Từ chối \\ \addlinespace
Cột sản phẩm &
  \rt{sale.order.line} &
  Sản phẩm của \textbf{dòng} trên đơn gốc --- mỗi phiếu đúng \textbf{một} sản phẩm &
  --- &
  \uat{10}/14 phiếu chưa gắn đơn gốc nên cột này trống \\ \addlinespace
Cột đơn gốc · chuyến &
  \rt{sale.order} · \rt{stock.picking.batch} &
  Mã đơn đặt hàng và mã chuyến giao lấy theo đơn đó &
  --- &
  \uat{4}/14 phiếu có đơn gốc \\ \addlinespace
Ảnh minh chứng &
  \rt{ir.attachment} &
  Số ảnh đã đính kèm phiếu &
  --- &
  \uat{1}/14 phiếu có ảnh (3 ảnh) \\
\end{hienbang}

\mloc

\begin{itemize}[leftmargin=*]
  \item \textbf{Trạng thái} --- 8 nhãn kể trên, mặc định \emph{``-- Tất cả trạng thái --''}.
        Giá trị lạ bị bỏ qua.
  \item \textbf{Từ khoá} --- rộng nhất trong portal: khớp một phần theo \textbf{mã yêu cầu},
        \textbf{mã đơn gốc}, \textbf{mã chuyến giao}, \textbf{tên sản phẩm} hoặc \textbf{mã
        sản phẩm}.
  \item \textbf{Khoảng ngày} --- lọc theo \emph{ngày tạo yêu cầu}, quy đổi đúng múi giờ tài
        khoản. Ngày gõ sai định dạng $\rightarrow$ băng báo ở đầu trang; ngày ngược
        $\rightarrow$ báo ngay tại thanh lọc và \textbf{không chạy truy vấn}.
  \item \textbf{Số phiếu mỗi trang} --- 10, 20, 50; mặc định \textbf{20}.
\end{itemize}

\mkiem{Tạo yêu cầu}

Không kiểm gì ở màn danh sách --- sang màn tạo yêu cầu, và màn đó kiểm lại toàn bộ từ đầu.

\mghi

\textbf{Không ghi gì.} Màn chỉ đọc.

\mhoi

\begin{ask}
\begin{enumerate}[leftmargin=*]
  \item \textbf{Danh sách gộp mọi cửa hàng mà không có bộ lọc cửa hàng}, cũng không có cột cửa
        hàng trên bảng. Người phụ trách nhiều cửa hàng khó biết phiếu nào của đâu.
  \item \textbf{Nhãn ``Nháp'' có trong bộ lọc} dù phiếu nháp không gửi được từ portal (xem màn
        sau).
  \item \uat{10}/14 phiếu trên UAT là dữ liệu mẫu nhập thẳng trong ERP, thiếu đơn gốc, sản phẩm
        và loại lỗi --- nếu cần bản chụp màn hình sạch để trình bày thì phải dọn trước.
\end{enumerate}
\end{ask}

% =============================================================================
\section{Màn Tạo yêu cầu đổi trả · bù hàng}
\rt{/portal/return/new}
\medskip

\mvao

Nút \emph{Tạo yêu cầu} ở màn danh sách. Cùng một đường dẫn phục vụ cả việc mở form (mở trang)
lẫn việc gửi form (bấm nút).

\mai

Phải đăng nhập và truy cập được ít nhất một cửa hàng; không thì bị đưa về danh sách kèm băng
báo chưa chọn cửa hàng. \textbf{Không chặn vai trò}.

\mhien

\begin{hienbang}
Ô \emph{Cửa hàng} &
  \rt{wujia.franchise.management} &
  \textbf{Mọi cửa hàng tôi truy cập được}. Chỉ có một cửa hàng thì ô này ẩn đi và gài sẵn. &
  --- &
  Administrator thấy \uat{2} cửa hàng \\ \addlinespace
Ô \emph{Loại lỗi} &
  \rt{wujia.return.issue.type} &
  Các loại lỗi \textbf{đang bật} (\rt{active = true}) &
  Theo số thứ tự cấu hình &
  \uat{5} loại: Hỏng bao bì · Sai sản phẩm · Hết hạn/cận date · Thiếu số lượng · Lỗi khác \\ \addlinespace
Ô \emph{Đơn hàng gốc} &
  \rt{sale.order} &
  Đơn \textbf{đã xác nhận} (\rt{state in (sale, done)}) của \textbf{mọi cửa hàng của tôi}, có
  \textbf{ngày đặt trong 10 ngày} trở lại. \textbf{Không} lọc lại theo cửa hàng vừa chọn ở ô
  bên cạnh. &
  Ngày đặt mới nhất trước &
  \uat{0} đơn tại ngày chụp --- đơn xác nhận gần nhất là 05/09, đã 15 ngày \\ \addlinespace
Ô \emph{Sản phẩm} &
  \rt{sale.order.line} &
  Các dòng hàng của đơn vừa chọn. Máy chủ dựng sẵn danh sách dòng hàng của \textbf{tất cả} đơn
  đủ điều kiện rồi nhúng vào trang, nên đổi đơn là đổi danh sách ngay, không phải gọi lại. &
  Theo thứ tự dòng trong đơn &
  --- \\ \addlinespace
Khối minh chứng &
  \rt{ir.attachment} &
  \textbf{3--5 ảnh} JPG/PNG, mỗi ảnh $\leq$ 5 MB; \textbf{tối đa 1 video} MP4/MOV
  $\leq$ 10 MB; tổng toàn bộ $\leq$ 30 MB &
  --- & --- \\
\end{hienbang}

\mloc

\begin{itemize}[leftmargin=*]
  \item \textbf{Số lượng yêu cầu} --- số thực lớn hơn 0 (khác màn Đặt hàng: ở đây \textbf{cho
        phép số lẻ}, không có bước nhảy, không có trần).
  \item \textbf{Thời gian mở hàng} --- \textbf{bắt buộc}, ô chọn ngày kèm giờ. Hệ thống nhận ba
        kiểu ghi ngày giờ; không đọc được thì báo lỗi và giữ lại những gì đã nhập.
  \item \textbf{Ngày sản xuất} --- không bắt buộc.
  \item \textbf{Ghi chú} --- cắt còn 5.000 ký tự, phần dư bị bỏ lặng lẽ.
  \item \textbf{Hai nút gửi khác nhau giữa hai nền tảng}: máy tính có \emph{Lưu nháp} và
        \emph{Gửi yêu cầu}; điện thoại \textbf{chỉ có} \emph{Gửi}.
\end{itemize}

\mkiem{Gửi yêu cầu}

\begin{kiembang}
1 & Truy cập được ít nhất một cửa hàng & Đưa về danh sách kèm băng báo chưa chọn cửa hàng \\
2 & Cửa hàng gửi lên là số và \textbf{nằm trong danh sách của tôi} &
    ``Cửa hàng không truy cập được.'' \\
3 & Có chọn đơn gốc và sản phẩm, cả hai là số &
    ``Vui lòng chọn đơn hàng gốc và sản phẩm.'' / ``Đơn hàng / sản phẩm không hợp lệ.'' \\
4 & Đơn đó \textbf{vẫn đủ điều kiện} --- kiểm lại ở máy chủ: đúng cửa hàng vừa chọn, đã xác
    nhận, trong 10 ngày &
    ``Đơn hàng không hợp lệ hoặc đã quá thời hạn 10 ngày.'' \\
5 & Dòng sản phẩm thuộc \textbf{đúng đơn đó} &
    ``Sản phẩm phải thuộc đơn hàng gốc của cửa hàng.'' \\
6 & Sản phẩm \textbf{đã được cấu hình chính sách bù hàng} (bật cờ bù, có đơn vị khai hao hụt,
    và tỷ lệ quy đổi hợp lệ) &
    ``Sản phẩm chưa được cấu hình chính sách bù hàng. Vui lòng liên hệ Ngô Gia.'' ---
    trên UAT \uat{4}/5 sản phẩm portal rơi vào nhánh này \\
7 & Loại lỗi còn bật & ``Vui lòng chọn loại lỗi.'' \\
8 & Số lượng lớn hơn 0 & ``Số lượng yêu cầu phải lớn hơn 0.'' \\
9 & Thời gian mở hàng đọc được & ``Vui lòng nhập thời gian mở hàng hợp lệ.'' \\
10 & Số ảnh trong khoảng \textbf{3--5} &
    ``Cần tải từ 3 đến 5 ảnh minh chứng.'' --- bấm \emph{Lưu nháp} thì \textbf{bỏ qua mức tối
    thiểu}, vẫn chặn mức tối đa \\
11 & Nhiều nhất 1 video & ``Chỉ được tải tối đa 1 video minh chứng.'' \\
12 & Mỗi tệp \textbf{đúng định dạng thật} và đúng dung lượng &
    ``Tệp không đúng định dạng hoặc vượt quá dung lượng cho phép.'' --- cùng một câu cho cả hai
    lý do, cố ý không cho biết đã sai cái nào \\
13 & Tổng dung lượng $\leq$ 30 MB & ``Tổng dung lượng minh chứng không được vượt quá 30 MB.'' \\
14 & Ghi được phiếu &
    ``Không thể gửi yêu cầu. Vui lòng kiểm tra lại thông tin và thử lại.'' ---
    \textbf{không} lộ chi tiết kỹ thuật ra màn \\
15 & Gắn được tệp vào phiếu &
    Phiếu vừa tạo \textbf{bị xoá hẳn}, quay lại form kèm lỗi --- không để lại phiếu rỗng \\
16 & (chỉ khi bấm Gửi) Đủ \textbf{3 ảnh} lúc chuyển sang \emph{Đã gửi} &
    Phiếu \textbf{bị xoá hẳn}, quay lại form \\
\end{kiembang}

\textbf{Về bước 12:} hệ thống đọc \textbf{nội dung thật} của tệp để biết đó là ảnh hay video,
không tin đuôi tệp và không tin loại tệp trình duyệt khai. Đổi tên một tệp bất kỳ thành
\rt{.jpg} không qua được bước này.

\mghi

\begin{itemize}[leftmargin=*]
  \item Tạo \textbf{một} bản ghi yêu cầu (\rt{wujia.return.request}), mã tự sinh theo dãy số
        (UAT đang ở \emph{RTN/26/00004}).
  \item Trạng thái ban đầu luôn là \textbf{Nháp}; bấm \emph{Gửi} thì hệ thống chuyển tiếp sang
        \textbf{Đã gửi} và ghi một dòng \emph{``Yêu cầu đã được gửi''} vào nhật ký trao đổi của
        phiếu.
  \item Người tạo = người đang đăng nhập; ngày tạo = lúc bấm nút. Cửa hàng ghi trên phiếu là
        cửa hàng chọn trong form, \textbf{không} phải cửa hàng đang chọn ở cookie.
  \item Đơn vị tính dùng để khai hao hụt lấy theo \textbf{cấu hình bù của sản phẩm}; sản phẩm
        chưa khai thì lùi về đơn vị ghi trên đơn gốc.
  \item Tạo các bản ghi tệp đính kèm, tách riêng nhóm ảnh và nhóm video.
  \item \textbf{Không} tạo đơn hàng, \textbf{không} động vào kho, \textbf{không} gửi email.
  \item Bất kỳ bước nào sau khi tạo mà hỏng thì phiếu bị \textbf{xoá hẳn} --- hệ thống không để
        lại phiếu mồ côi hay tệp mồ côi.
\end{itemize}

\mhoi

\begin{ask}
\begin{enumerate}[leftmargin=*]
  \item \textbf{Cửa sổ 10 ngày đang làm màn này không dùng được.} Mốc tính từ \emph{ngày đặt}
        của đơn đã xác nhận, đếm \textbf{theo giờ} chứ không theo ngày. Trên UAT có \uat{14} đơn
        đã xác nhận nhưng \uat{0} đơn lọt cửa sổ (đơn gần nhất 05/09, cách 15 ngày)
        $\Rightarrow$ \textbf{không ai tạo được yêu cầu nào}, ô chọn đơn trống kèm dòng đỏ.
        Anh chốt giúp: giữ 10 ngày hay nới ra, và tính từ ngày đặt hay ngày \textbf{giao hàng}?
  \item \textbf{Ô chọn đơn không lọc theo cửa hàng ở ô bên cạnh.} Người nhiều cửa hàng chọn
        lệch cặp cửa hàng--đơn sẽ nhận câu ``Đơn hàng không hợp lệ hoặc đã quá thời hạn 10
        ngày'' --- câu này \textbf{nói sai nguyên nhân}.
  \item \textbf{Phiếu nháp là ngõ cụt.} Máy tính có nút \emph{Lưu nháp}, nhưng portal không có
        chỗ nào gửi phiếu nháp đã lưu: màn chi tiết chỉ đọc. UAT đang có \uat{3} phiếu kẹt.
        Anh muốn (a) bỏ nút Lưu nháp, hay (b) thêm nút \emph{Gửi} ở màn chi tiết?
  \item \textbf{Điện thoại không có nút Lưu nháp} --- hai nền tảng khác luật cho cùng một việc.
  \item \textbf{Sản phẩm chưa cấu hình bù thì chặn cả đổi và trả}, dù hai phương án đó không
        cần chính sách bù. Trên UAT \uat{4}/5 sản phẩm đang hiện trên màn Đặt hàng rơi vào đây
        --- cửa hàng đặt được nhưng không khiếu nại được.
  \item \textbf{Mỗi phiếu đúng một sản phẩm.} Một đơn hỏng ba mặt hàng phải làm ba phiếu, mỗi
        phiếu tự chụp đủ 3 ảnh. Anh xác nhận giữ nguyên.
  \item \textbf{Video tải lên không bao giờ hiện lại cho cửa hàng.} Hệ thống nhận, kiểm và lưu
        video, nhân sự Ngô Gia xem được trong ERP, nhưng màn chi tiết trên portal chỉ hiện ảnh.
        Cửa hàng không có cách nào xem lại video mình vừa gửi.
  \item \textbf{Cửa hàng không sửa và không huỷ được phiếu sau khi gửi} --- portal không có nút
        nào cho việc đó.
\end{enumerate}
\end{ask}

% =============================================================================
\section{Màn Chi tiết yêu cầu}
\rt{/portal/return/<int:request_id>}
\medskip

\mvao

Bấm một phiếu ở màn danh sách; hoặc được đưa tới ngay sau khi tạo yêu cầu thành công (kèm băng
báo vừa tạo xong).

\mai

Phiếu phải thuộc \textbf{một trong các cửa hàng của tôi}. Không thoả thì bị đưa về danh sách
kèm băng \emph{không tìm thấy} --- \textbf{cùng một câu} cho cả trường hợp phiếu không tồn tại
lẫn trường hợp phiếu của cửa hàng khác, cố ý không cho dò số phiếu.

\mhien

\begin{hienbang}
Đầu phiếu &
  \rt{wujia.return.request} &
  Mã yêu cầu, ngày tạo, cửa hàng, huy hiệu trạng thái (tám nhãn như màn danh sách) &
  --- & --- \\ \addlinespace
Khối nội dung &
  \rt{wujia.return.request} &
  Loại lỗi, thời gian mở hàng, ngày sản xuất, số lượng yêu cầu, ghi chú của cửa hàng &
  --- & --- \\ \addlinespace
Khối minh chứng &
  \rt{ir.attachment} &
  \textbf{Chỉ ảnh}; mỗi ảnh mở bằng đường dẫn tải riêng của phiếu. Video \textbf{không} hiện. &
  --- &
  \uat{1}/14 phiếu có ảnh \\ \addlinespace
Khối phản hồi &
  \rt{reject_reason} · \rt{approval_note} &
  \textbf{Lý do từ chối} khi bị từ chối, hoặc \textbf{ghi chú duyệt} khi đã duyệt. Ghi chú nội
  bộ của Ngô Gia \textbf{không} lộ ra portal. Chưa có gì thì ghi ``chưa có phản hồi''. &
  --- & --- \\ \addlinespace
Khối tiến độ bù &
  \rt{wujia.compensation.allocation} &
  Chỉ hiện khi Ngô Gia \textbf{đã chốt phương án}. Với phương án \emph{Bù hàng}: số lượng
  duyệt, đã lên đơn, đã bù, còn thiếu, thanh phần trăm (\textbf{chặn trần 100\%}), và danh sách
  đơn bù kèm nhãn giao hàng. &
  --- &
  \uat{0} bản ghi phân bổ $\Rightarrow$ khối luôn trống, kể cả \uat{3} phiếu đã chốt phương án
  bù \\ \addlinespace
Nhãn giao của đơn bù &
  \rt{stock.picking} &
  Năm nhãn: \emph{Chưa tạo phiếu giao} · \emph{Chưa giao} · \emph{Đã giao n/m phiếu} ·
  \emph{Đã giao đủ} · \emph{Đơn bù đã bị hủy} &
  --- & --- \\
\end{hienbang}

Khi \textbf{mọi đơn bù đều đã bị huỷ}, màn hiện thêm một cảnh báo riêng --- để cửa hàng biết
quyền lợi đã đóng lại và phải tạo yêu cầu mới, thay vì ngồi chờ hàng.

\mloc

Không có ô lọc hay ô nhập nào.

\mkiem{bất kỳ nút nào}

Màn \textbf{chỉ đọc} --- không có nút nào ghi dữ liệu. Duy nhất các ảnh minh chứng là liên kết
mở sang đường dẫn tải (màn sau).

\mghi

\textbf{Không ghi gì.}

\mhoi

\begin{ask}
\begin{enumerate}[leftmargin=*]
  \item \textbf{Cửa hàng không có nút nào ở màn này} --- không gửi được phiếu nháp, không huỷ,
        không bổ sung ảnh, không trả lời lại Ngô Gia. Anh muốn thêm gì trước khi chạy thật?
  \item \textbf{Thanh tiến độ bù chưa chạy được trên UAT} vì chưa có bản ghi phân bổ nào. Cần
        dựng một ca mẫu (duyệt $\rightarrow$ lên đơn bù $\rightarrow$ giao một phần) để nghiệm
        thu đúng nhãn \emph{Đang bù một phần}.
  \item \textbf{Thanh phần trăm bị chặn trần 100\%} --- bù dư hơn số duyệt vẫn hiện đúng 100\%.
\end{enumerate}
\end{ask}

% =============================================================================
\section{Xem · tải minh chứng}
\rt{/portal/return/<int:request_id>/attachment/<int:att_id>}
\medskip

\mvao

Không phải màn --- đây là đường dẫn trả về chính tệp ảnh, gọi ngầm khi màn chi tiết hiện ảnh
thu nhỏ, và mở ra tab mới khi người dùng bấm vào ảnh.

\mai

Ba lớp kiểm, theo thứ tự:

\begin{kiembang}
1 & Truy cập được ít nhất một cửa hàng & Từ chối truy cập \\
2 & Phiếu tồn tại và thuộc \textbf{cửa hàng của tôi} & Báo không tìm thấy \\
3 & Tệp \textbf{thuộc đúng phiếu đó} (nằm trong nhóm ảnh hoặc nhóm video của phiếu) &
    Từ chối truy cập --- \textbf{không} cho mượn đường dẫn này để đọc tệp bất kỳ trong hệ thống \\
4 & Tệp còn tồn tại & Báo không tìm thấy \\
\end{kiembang}

\mhien

Trả về đúng nội dung tệp, đặt chế độ \textbf{xem thẳng trong trình duyệt} chứ không ép tải
xuống. Video được phép đi qua đường dẫn này nhưng \textbf{không màn nào trỏ tới} --- xem mục 7
màn Tạo yêu cầu.

\mloc

Không có.

\mkiem{một ảnh}

Xem bảng ba lớp ở trên.

\mghi

\textbf{Không ghi gì.}

\mhoi

\begin{ask}
Ảnh minh chứng đi qua đường dẫn có kiểm quyền nên \textbf{không} đọc được nếu chỉ có số tệp.
Không có điểm nào cần BA chốt ở đường dẫn này.
\end{ask}

% =============================================================================
\section{Đường dẫn cũ của phân hệ Đổi trả}
\rt{/portal/return-request-list}
\medskip

Đường dẫn từ bản v14, được giữ lại để chuyển hướng \textbf{vĩnh viễn} sang \rt{/portal/return}.
Khoảng 1.500 tài khoản nhượng quyền có thể còn dấu trang hoặc email cũ trỏ vào đó.

Đường dẫn này \textbf{không đòi đăng nhập} --- ai mở cũng được chuyển hướng, rồi mới bị màn
đích đòi đăng nhập. Đây là cố ý: chuyển hướng trước, hỏi đăng nhập sau, để người dùng không
mất đường dẫn đích sau khi đăng nhập xong.

\begin{ask}
Không có điểm cần chốt --- nhưng nếu BA muốn khai tử đường dẫn cũ sau một mốc thời gian nào đó
thì cho biết mốc, dev sẽ gỡ.
\end{ask}
